\section{Abstract}\label{abstract}

\section{Introduction}\label{introduction}

Our experience of the world can be described as a continuous stream of
sensory information that flows from the external environment into our
brains through our sensory systems. However, empirical studies suggest
that our memories for such experiences are organized according to
contextual features such as perceptual, semantic, temporal and spatial
similarities among the elements of an experience. How do our brains
distill our rich, continuous perceptual experience into meaningful and
manageable units in memory? This question is complex and difficult to
approach because this process likely requires a coordinated interaction
between perception, attention, working memory and long-term memory
processes. While a great deal has been revealed about the neural systems
and processes that contribute to memory formation for individual items
and their associated contexts (refs), the cognitive and neural
mechanisms associated with parsing experience into `mnemonic events' has
yet to be fully explored. Event segmentation theory (EST), a
computational/theoretical model of how the brain parses ongoing
experience into events, posits that during perception the human brain
maintains an `event model'. An event model can be defined as a working
memory representation of `what is happening now' (Kurby and Zacks).
Event models are thought to contain contextually relevant details of an
ongoing experience in order to generate predictions about future sensory
input, which is an advantageous to an organism in a dangerous and
uncertain world. When there is a substantive change in the perceptual
experience (for example, walking through a door way out of one room and
into another) or predictability breaks down, EST suggests that the
previous event model is abandoned (or perhaps, flushed from working
memory) and a new event model is generated. This process (known as
`event segmentation') is thought to help organize our experience into
meaningful units by integrating information that is perceived to be in
the same event, while keeping across-event information separate.

Emerging evidence suggests that one consequence of event segmentation is
that associative memory for information encoded in distinct events is
reduced relative to information encountered in the same event (Ezzyat et
al 2011, Dubrow et al 2013, Heusser et al 20XX). In these studies,
events are operationalized as temporally adjacent items with overlapping
contextual features. Ezzyat at al (2011) found that cued recall
performance is reduced for pairs of items that spanned a contextual
shift (or an `event boundary'), and temporal order memory for
across-boundary information is also reduced (DuBrow \& Davachi, 2013;
Heusser et al, 20XX). According to event segmentation theory, this may
occur because the online maintenance of an event model facilitates the
integration of within-event information, and event boundaries disrupt
integration processes. Bolstered by predictions from EST, these results
suggest the existence of a neural mechanism whereby items are maintained
and integrated within an event, but not across separate events.

In a separate literature, a number of studies have provided evidence
that sustained and synchronized neural activity (i.e.~neural
oscillations) support the maintenance of information in working memory
(citations). Low frequency oscillatory power (i.e.~Theta: 4-8 Hz) during
a working memory delay period has been shown to increase with the number
of items in working memory (citations), and the duration of the increase
scales with the length of the delay period (citations). Interestingly,
recent evidence suggests that maintenance-related increases in theta
power appear to be specifically related to working memory for temporal
relationships among items, as opposed to the items themselves (frank and
charan). Thus, theta power increases appear not only to track the amount
of information in working memory, but seems to be most strongly
associated with the maintenance of the associative relationships among
items.

Theta oscillations have also been linked to the successful long-term
memory encoding (citations) as well as retrieval (citations). Numerous
studies have found increased theta power for items that are subsequently
remembered as compared to items that are later forgotten (citations, but
see Long et al). Furthermore, a subset of these studies suggest that
theta oscillations are involved in the binding between an item and its
context (i.e.~a word seen in a particular color), rather than indexing
memory for an item alone (Summerfield; Hanslmayr).). Taken together,
these data suggest that memory-related theta oscillations may play a
more general role in mnemonic binding or integration, whether it's an
item to its context or an item to another item.

The relationship between lower frequency oscillations and memory has
been studied extensively. Perhaps not surprisingly, most other frequency
bands have also been linked to mnemonic processes in various
experimental designs (citations). The frequency band that has received
probably the least amount of attention with respect to memory is the
beta (13-30 Hz) band. However, there is also some theoretical work
suggesting that beta oscillations may play a role in the maintenance of
cognitive states (for review, see Engel and Fries). While the majority
of this evidence comes from studies involving the maintenance of motor
postures (i.e.~holding an arm in a particular posture) and/or motor
plans, a few studies suggest beta may be important for the maintenance
of mnemonic information (talon bauldtry papers) as well as the
maintenance of an attentional state (buschman and miller). One recent
study found that when recording from the prefrontal cortex of monkeys
performing a 2-item order maintenance task, beta oscillations increased
during the delay period and beta phase was related to subsequent
temporal order memory success. Thus, beta oscillations may also be
involved mnemonic processes and play a particular role in the
maintenance of a cognitive state or `event model' (Engel and Fries).

In this study, we utilized magnetoencephalography (MEG) to ask whether
dynamics in oscillatory power might support the formation of episodic
events. Previous electrophysiological studies of memory have focused on
within-trial spectral and temporal dynamics to examine the time course
of memory-related oscillatory processes. However, relatively fewer
studies have considered the modulation of oscillatory power at longer
timescales (\textgreater{} 5 seconds). In this study, we focused our
analyses on across-trial changes in power to investigate how slower
fluctuations in oscillatory power are modulated by the event structure
in our paradigm and are related to the successful encoding of episodic
events. Guided by models of event segmentation and strong evidence in
support of a relationship between neural oscillations and memory, we
hypothesized that oscillatory power in the theta and beta band would
increase within an event as more information is accumulated and
integrated into the event model. At event boundaries, we predicted that
power would drop significantly as within-event representations might be
flushed from working memory. Finally, we hypothesized that within-event
power accumulation would be related to temporal order memory on a later
test. Specifically, we predicted that successful within-event
integration would be associated with a monotonic increase in power
within an event that might index accumulating within-event information.

\section{Methods}\label{methods}

\subsubsection{Participants}\label{participants}

Twenty healthy native English speakers (XX males, age range: XX) were
recruited from New York University and the greater New York Metropolitan
area participated in the MEG experiment. All participants gave informed
written consent in accordance with the University Committee on
Activities Involving Human Subjects (UCAIHS). We excluded one subject
for chance performance on the order memory test and one subject did not
complete the study due to drowsiness. We excluded these participants,
leaving 18 subjects for the remaining MEG analyses.

\subsubsection{Materials}\label{materials}

Details for this experiment are published in Heusser et al., 20XX. The
experimental design is identical to Experiment 2 in that manuscript.
Briefly, stimuli consisted of 576 gray-scale pictures of objects
collected from various online sources. Colored borders for the objects
were generated by selecting 24 unique colors from a color continuum
ranging from {[}0,0,0{]} to {[}255,255,255{]} RGB values. Colors were
chosen to be maximally distinct from each other. Within the experiment,
the order of the stimuli varied for each subject and color-stimulus
pairings were randomized across subject.

\subsubsection{Encoding}\label{encoding}

During encoding, participants made pleasant/unpleasant decisions on
trial-unique objects that were paired with a colored background frame
(see Figure 1). Specifically, participants were instructed to imagine
each object in the color of the background frame and press a button to
indicate whether or not the combination was pleasant. Additionally, we
instructed participants to integrate the objects together by imagining
them interacting with each other. We instructed participants to use this
encoding strategy because it would help them with a later temporal order
memory test. The background color frame remained the same for 6
consecutive trials and then switched to a new color. We operationalized
objects that were studied consecutively in the same color as an `event'
and switches in color as an `event boundary'. There were 6 events
(totaling 36 objects) in each encoding block and 16 encoding/test blocks
across the experiment. Each object was on the screen for 2.5 seconds,
followed by a 2 second inter-trial interval and a .5 second fixation
period. The timing of the task was fixed (i.e.~not jittered).

\subsubsection{Color Memory Test}\label{color-memory-test}

After each encoding block, we tested two types of memory. First, we
tested object-color associative memory. In this test, a previously
studied object appeared in the middle of the screen and two colors
appeared below it. One of the colors was studied with the object (target
color), while the other was a lure color studied in a neighboring event.
Subjects were instructed to indicate which of the two colors was studied
with the object and rate their confidence. Thus, there were four
possible responses during the test: High confidence target, low
confidence target, high confidence lure and low confidence lure. There
were two conditions in the color memory test. The `boundary' condition
consisted of trials that were studied concurrent with a color switch.
The `non-boundary' condition was comprised of trials from the middle of
an event (specifically in the fourth position of an event). The test was
self-paced with a mandatory .5 second fixation period between each test
trial.

\subsubsection{Temporal Order Memory
Test}\label{temporal-order-memory-test}

After the color memory test, there was a temporal order memory test. In
this test, two previously studied objects were presented side by side
without the colored background frame. Participants were asked to
indicate which of the two objects appeared first in the sequence. Like
the color memory test, there were four possible responses (see above).
In the `within-event' condition, objects from the 2nd and 6th
within-event position were presented together and participants were
asked to indicate which came first (position two would be the correct
answer). In the `across-boundary' condition, objects studied in the 5th
position of an event and the 3rd position of the next event were
presented together (position 5 would be the correct answer). Critically,
all tested object pairs were separated by 3 intervening trials during
encoding. Thus, while objects in the `within-event' condition came from
the same event, objects in the `across-boundary' condition were studied
in neighboring events and the amount of elapsed encoding time between
the two test items for the two conditions was identical.

\subsubsection{MEG recordings and data
processing}\label{meg-recordings-and-data-processing}

MEG data were recorded using a 157-channel whole-head axial gradiometer
system (KIT, Kanazawa, Japan). Three reference channels seated above the
MEG system were also recorded and used to remove ambient electromagnetic
environmental noise from the data. MEG data were acquired with a
sampling rate of 1000 Hz, which was filtered online between XX and XX
Hz, with a notch filter at XX Hz. To measure head position, 5
electromagnetic coils were attached to a participant's head. Coil
locations were determined by registering scalp coil positions with 3D
digitized head shape data (software: Source Signal Imaging, Inc.;
hardware: Polhemus, Inc.), which was collected before MEG recording. The
anatomical locations used to register the coils with the head shape data
were the nasion and the left and right periauricular points. The coils
were localized to the MEG sensors at the beginning and end of the
experiment.

MEG data were preprocessed in the following way: raw data were loaded
into Matlab (version 7.10, Mathworks) and malfunctioning channels
(average per subject: XX) were immediately removed and interpolated
using a nearest neighbors approach (CITATION). Then, data were denoised
using a time-shifted principal components analysis approach, which
removed ambient environmental noise using 3 reference channels
(CITATION). The remaining preprocessing steps used the Fieldtrip M/EEG
software package (CITATION). The data were band-pass filtered from 1-100
Hz. Data were epoched from -4 to 4 seconds surrounding trial onset to
assure adequate time for spectral estimation of both pre- and
post-stimulus activity. The data were down-sampled to 500 Hz to speed
processing time in later steps. Finally, to facilitate interpretation of
topographies, we transformed the MEG data from axial to planar gradient.
One advantage of this linear transformation is that planar signal
amplitude is typically largest directly above the source, whereas axial
signal amplitudes are typically maximal on either side of the neural
source of the signal.

After preprocessing, the MEG data were examined for artifacts. The
artifact rejection approach we took was three-fold: First, excessively
noisy trials and channels were removed using Fieldtrip's visual artifact
rejection `trial summary' feature. Specifically, channels trials for
which the across-trial variance exceeded 3 standard deviations from the
mean were removed from the analysis. Then, independent components
analysis was implemented to remove components related to eye blinks, eye
movements, and heartbeat-related artifacts. Finally, remaining trials
were visually inspected and epochs containing any remaining artifacts
were removed from the dataset (\textasciitilde{}XX\% of trials).

\subsubsection{Time-Frequency Analyses}\label{time-frequency-analyses}

A time-frequency transformation was applied to each epoch (-4 to 4
seconds, 50 ms sliding window, zero-padded) using a morlet wavelet
approach (number of cycles=6) in Fieldtrip, estimating spectral power
from 1 to 100 Hz in steps of 1Hz. This analysis resulted in
time-frequency spectrograms representing oscillatory power in each
time-frequency-sensor point for each trial and each subject. This
relatively long epoch window allowed us to analyze data during the
`stimulus on' period (0 to 2.5 seconds) as well as the inter-trial
periods (-2.5 to 0 seconds) while avoiding edge artifacts particularly
in the low frequencies. For analyses where power was averaged across the
entire trial, data from the trial N (0 to 2.5 seconds) and the
pre-stimulus period of the next trial (N+1; -2.5 to 0 seconds) were
averaged. This was because we centered the time-frequency analyses on
the onset of each trial.

\subsubsection{Power Accumulation
Analyses}\label{power-accumulation-analyses}

To determine whether oscillatory power increased as a function of within
event position, we adopted a general linear model approach (GLM). For
each trial (\textasciitilde{}XX after artifact rejection), we averaged
oscillatory power across the entire trial (0 to 5 seconds) for each
frequency (1-100 Hz) and each sensor (n=157). We then set up our `power
accumulation' model. The power accumulation model was simply a
mean-centered vector that linearly increased from 1 to 6 within each
event and reset back to 1 at event boundaries (see Figure 2). Put
another way, the accumulation model indexed the within-event position of
each trial. Also included in the GLM was a regressor modeling the mean
signal across the entire experiment as well as a regressor to account
for mean signal differences across the 16 blocks. We then regressed the
power data (separately for each band and sensor) against this model to
obtain a t-statistic representing the fit of the oscillatory power data
to the `accumulation' model. A model fit statistic (t-value) was
obtained for each frequency bin and subject (as well as each sensor and
time point depending upon the analysis). Finally, to look for
reliability of these model fits across subjects, we performed a
group-level one-sample t-test at each frequency bin-sensor data point.

\subsubsection{Source Localization
Analysis}\label{source-localization-analysis}

Beamformer analysis (DICS; CITATION) was performed to estimate neural
sources of the power accumulation effects. Each subject's data was first
registered to a structural template MPRAGE brain (FSL, XX) by aligning
anatomical landmarks (nasion, left/right periauricular points) from the
digitized head shape data to the structural brain image. Then,
oscillatory power (for the bands of interest: delta, theta and beta) was
computed in source space for each trial during the stimulus on period
(0-2.5 seconds). We used only the `stimulus on' periods in this analysis
since this time period was the most significant in the sensor space
analyses (see Figure 3). The source-space data were fit to the same
power accumulation model described above, to identify putative
generators of the ramping within-event oscillatory neural activity.
Group-level reliability was then calculated with a one-sample t-test
across subjects.

\section{Results}\label{results}

\subsubsection{Encoding}\label{encoding-1}

On each encoding trial, participants were instructed to imagine an
object in the color of the background frame and decide if the
combination was pleasant or unpleasant (62\% rated as pleasant).
Encoding task response times (RTs) varied as a function of within-event
position {[}F(5,80)=16.44, p\textless{}.0001{]}, where follow-up
pairwise t-tests reveal that RTs to boundary items were significantly
slower than RTs to other positions {[}2-5, all p's\textless{}.001,
corrected for multiple comparisons{]} and there were no other
differences. This boundary-related RT cost replicates a previous study
using this behavioral paradigm (Heusser et al, 20XX).

\subsubsection{Temporal Order Memory}\label{temporal-order-memory}

After each encoding block, participants were asked to decide which of
two previously studied items occurred earlier in the sequence (see
Methods for more details). We predicted that within-event temporal order
memory would be greater than across-event order memory. Consistent with
our hypothesis, temporal order memory was significantly more accurate in
the within-event condition than the across-boundary condition
{[}within-event condition: Mean=.73, SD=.12; across-boundary condition:
Mean=.57, SD=.09; t(17)=5.56, p\textless{}.0001{]}. Critically, the
temporal lag between the retrieval probes in both conditions was fixed
at three and thus, the only difference between conditions was the
occurrence of an intervening event boundary in the across-boundary
condition. These data suggest that event boundaries result in reduced
temporal order memory (relative to within-event temporal order memory),
consistent with previous studies (Heusser et al, 20XX; Dubrow et al,
2013). Next, we tested whether response times during the temporal order
retrieval task varied by condition. We hypothesized that even when
participants make a correct temporal order response, they might be
slower to respond to across-boundary order test trials as compared to
the within-event test trials. Consistent with this prediction
across-boundary test trial RTs {[}M=3.79s, SD= 1.67{]} were
significantly slower than within-event test trial RTs {[}M=3.43s,
SD=1.22; t(17)=2.36, p\textless{}.05{]}. Thus, in addition to reduced
accuracy, when participants made a correct response in the
across-boundary order condition, they were significantly slower.

\subsubsection{Object-Color Memory}\label{object-color-memory}

We hypothesized that object-color memory at event boundaries would be
facilitated relative to object-color memory for trials in the middle of
an event (specifically, in the 4th position of an event). To test this
hypothesis, we probed memory for the object-color association for these
two positions. After each encoding block, we presented a previously
studied object along with a target color (i.e.~that color that was
studied with the word) and a lure color (a color from a neighboring
event). Consistent with our hypothesis, object-color associative memory
was significantly greater at boundaries {[}M=87.27\%, SD=8.41\%{]} then
at non-boundaries {[}M=79.08\%, SD=9.25\%{]}, suggesting that event
boundaries facilitated object-color binding {[}t(17)=5.54,
p\textless{}.001{]}. Next, we asked whether correct responses were
speeded in the boundary condition, relative to the non-boundary
condition. We hypothesized that even when conditions were matched for
accuracy, boundary trials would be accessed more quickly than
non-boundary trials. Indeed, correct boundary RTs {[}M=2.05s, SD=.58{]}
were on average half of a second faster than correct non-boundary RTs
{[}M=2.54s, SD=.84{]}, suggesting that color-object memory on boundary
trials was accessed significantly faster than non-boundary trials
{[}t(17)=-4.85, p\textless{}.001{]}. Both the temporal order and color
memory findings are consistent with previous studies in our lab (Heusser
et al 20XX). The remainder of this manuscript will focus on oscillatory
activity during the encoding of events related to the integration of
within-event information.

\subsubsection{Oscillatory power increases within
events}\label{oscillatory-power-increases-within-events}

To test our hypothesis that oscillatory power would increase within an
event (i.e.~6 objects studied in with the same color) and drop off at
event boundaries, we adopted a model-based regression approach (Figure
2, see Methods for details). For each subject, we fit an `accumulation'
model to the trial-level oscillatory power estimates. As described in
the Methods, the accumulation model was simply a function that increased
linearly within each event (i.e.~6 adjacent items with the same color
backgrund; see Figure 2 for visual depiction). Power was averaged within
each trial (0 - 2.5 seconds) and the modulation of power across trials
was fit to the accumulation model. First, we performed this across-trial
regression analysis for each frequency band separately (1-100 Hz),
averaging across all MEG sensors. The results show well-defined peaks in
the distribution of model fits in delta (2 Hz) and theta (7 Hz), as well
as two peaks in the beta band (16 and 26 Hz), indicating that these
frequency bands significantly fit the accumulation model
{[}p's\textless{}.001, Figure 2{]}. In addition to the parametric
statistics, we also performed a permutation analysis, where for each
iteration (n=1000), we shuffled the trial-level data (within-subject)
and refit the model for every subject and then recomputed the
group-level statistics. The result was a null distribution of
group-level model fit significance values for each frequency band. The
non-parametric statistics were very similar to the parametric tests and
each frequency peak (in delta, theta and beta) remained significant
(p\textless{}.001). As an additional control analysis, we recomputed the
fits between the power data in each frequency but shifted the phase of
the accumulation model. We predicted that the model that was `in-phase'
with the event structure of our task would be the most significant,
while `out-of-phase' models would not show a reliable relationship to
the data. There were 5 possible phase-shifted models. None of the 5
out-of-phase models reached significance in any frequency band at our
statistical threshold (p\textgreater{}.001). Thus, the model that was
in-phase with the event structure of the task best fit the data. In
summary, power in delta, theta and beta significantly co-varied with the
event structure of the experiment, suggesting that these bands may play
a role in event formation and segmentation.

\subsubsection{Temporal specificity of power accumulation
effects}\label{temporal-specificity-of-power-accumulation-effects}

Our next analysis aimed to characterize what time periods during each
trial showed the most robust power accumulation effects. One possibility
is that the stimulus-evoked response increase in magnitude across the
event and drive the significant model fits reported above.
Alternatively, if power was gradually increasing across the event, one
might expect the linear increases in power to be present for the
duration of the trial and into the inter-trial interval. For each trial,
we computed a time-frequency spectrogram (-1 to 4 seconds in 50ms bins,
1-100 Hz in 1 Hz bins) and then fit the accumulation model across trials
to each time-frequency data point. This analysis resulted in a single
time-frequency spectrogram (averaged across all MEG sensors), where each
data point represents the fit to the accumulation model at a particular
time point in a particular frequency (Figure 2). The result of this
analysis suggests that model fits were maximal during the time that the
stimulus was on the screen, but were also robust in the post-stimulus
time window (after 2.5 seconds when the stimulus was removed from the
screen), suggesting that the power accumulation effects were not driven
by the stimulus-evoked component and not limited to the temporal
interval when the stimulus was on the screen. Thus, these effect do not
appear to be driven by stimulus-evoked activity alone, but rather more
likely reflects cognitive processes evolving at longer timescales
presumably tracking the event structure of the task.

\subsubsection{Spatial specificity of power accumulation
effects}\label{spatial-specificity-of-power-accumulation-effects}

Next, we sought to characterize the topographic distributions of the
within-event power accumulation effects. We limited our analyses to
frequency bands that displayed peaks in the distribution of accumulation
model fits (specifically, the delta, theta and beta bands). We found
that both delta and theta oscillatory accumulation model fits were
significant and maximal over posterior sensors {[}Figure 3,
p\textless{}.001, cluster threshold: p\textless{}.001, corrected at
p\textless{}.05{]}. On the other hand, model fits to beta power were
significant in more sensors distributed across most of the scalp
{[}p\textless{}.001, cluster threshold: p\textless{}.001, corrected at
p\textless{}.05{]}.\\
To verify that the within-event power increases that we've characterized
thus far were driven by gradual increases in power within an event that
drops at event boundaries (rather than being driven by another pattern
of activity or alternatively, an artifact of the way we modeled the
data), we extracted power estimates from sensors that showed a
significant fit to the accumulation model. Then, separately for each
frequency band of interest (delta, theta, beta), we averaged the power
data over significant sensors. This resulted in a single time course of
oscillatory power for each frequency band across the entire duration of
the experiment. To increase our signal to noise ratio, we downsampled
our time courses, binning the power time courses into 500 millisecond
chunks and and averaging. We also averaged over all 16 encoding blocks
(and across subjects). Plotted in Figure 3 is the time course of
oscillatory power for delta, theta and beta across an entire encoding
block (i.e.~36 trials / 6 events). Each gray line represents the slope
of the power increase within an event. Visual inspection of the time
courses clearly display accumulating power within an event that drops at
the event boundaries in all three frequency bands of interest. This
analysis provides compelling evidence that oscillatory power in the
delta, theta and beta bands tracks the event structure in our task.

\subsubsection{Ruling out potential
confounds}\label{ruling-out-potential-confounds}

We performed a series of control analyses to consider possible confounds
and alternative explanations. First, we noted that oscillatory power
appeared to increase not only within an event, but also across the
duration of a study block (see Figure 3). In other words, we observed
power increases within an event but also power increases over a block of
36 trials. To test whether our accumulation model analyses were driven
by power changes across an encoding block (rather than power increases
within an event), we regressed out block-level linear trends and
recomputed the model fits on the residuals of this regression. This did
not account for the results, as the effects remained statistically
robust (see Supplemental Figure XX). We also noticed that overall power
in the first event of a block was notably lower than the rest of the
events. Thus, we removed the first event from the analyses and
recomputed the accumulation model analyses. Again, this did not change
the pattern of results (see Supplemental Figure XX). Another possibility
is that evoked responses (i.e.~ERPs) might contribute to this
accumulation effect, rather than induced oscillatory activity. To rule
out this possibility, we subtracted out the average evoked response from
each trial and recomputed the accumulation analysis. This did not change
the results, suggesting that trial-evoked responses did not explain the
power accumulation effects. Finally, we tested whether the accumulation
effects were sensitive to baseline correction. We baseline corrected
each trial to the pre-stimulus period (-1000 to -500 ms). Interestingly,
the effects were eliminated (See Supplemental Figure XX), which is
consistent with the idea that the power accumulation effects were not
limited to the periods when the stimulus was on the screen, but may also
exist in the pre-stimulus periods. This interpretation is also supported
by the fact that the pre-stimulus oscillatory power in delta, theta and
beta fits the accumulation model (See Figure 3). Together, these data
provide compelling evidence that oscillatory power in delta, theta and
beta increases within an event and drops at event boundaries, and that
this dynamic pattern likely reflects cognitive processes that evolve
over the timescale of a temporally-extended episodic experience.

\subsubsection{Source Localization of the power accumulation
effects}\label{source-localization-of-the-power-accumulation-effects}

The purpose of the next analysis was to 1) localize the source(s) of
within-event power accumulation signals and 2) test whether different
frequency bands localized to different brain areas. To achieve this, we
utilized a beamforming procedure (see Methods for details) to estimate
oscillatory source power for each trial in three frequency bands of
interest (delta, theta and beta). Given that we observed the strongest
accumulation effects during the `stimulus on' period in the trial (0 to
2.5 seconds), we averaged power across this epoch for each trial. We
then fit the source space power data to the same power accumulation
model as described above. We (primarily) focused our analyses on the
peaks in the frequency distribution of accumulation model fits (as shown
in Figure 2), since these were the frequency bands that best fit our
power accumulation model. Delta (2 Hz) power accumulation was localized
to the right hippocampus and right cerebellum. High theta (7 Hz) was
localized to the right paracentral lobule. Low beta (16 Hz) localized to
the right MTL, right inferior frontal junction, and left premotor
cortex. High beta (26 Hz) localized to the cerebellum bilaterally, left
lateral occipital complex (LO) and right parietal cortex. Finally, given
recent evidence suggesting that low theta (\textasciitilde{}3-4 Hz) is
related to episodic memory in humans (Watrous et al; Jacobs et al), we
source localized low theta (4 Hz) power accumulation. This analysis
localized low theta power accumulation to the left hippocampus/MTL, the
left cerebellum (vermis VIIIa), and the left superior frontal sulcus.

Visual inspection revealed overlap in the source localization for theta
(4 Hz) and low beta (16 Hz). To quantify this overlap, we ran a
conjunction analysis (joint p\textless{}.0001, individuals maps
thresholded at p\textless{}.01). This analysis revealed that source
localization for beta and theta overlapped in the left MTL and a left
insular region (Figure 3). Thus, while low frequency power accumulation
(4 Hz and lower) was localized primarily to the hippocampus and MTL,
higher frequency beta power accumulation was localized to frontal and
parietal regions (as well as MTL and LO) and there was overlap between
the theta and beta source localization in the left MTL and insula.

\subsubsection{Power accumulation and subsequent temporal order
memory}\label{power-accumulation-and-subsequent-temporal-order-memory}

In the next set of analyses, we asked whether increasing power within an
event was related to successful temporal order memory. We focused our
analyses on sensors that showed a significant fit to the accumulation
model in our three bands of interest (delta, theta and beta). For each
trial, we extracted power from sensors that displayed significant
accumulation model fits (averaging across sensors). Then, we binned
trials by the within-event trial position (1-6) and also by later
within-event temporal order memory. To test for differences in the
pattern of power accumulation between hits and misses, we performed
repeated-measures ANOVAs testing for a significant linear trend within
an event. We excluded boundary trials (i.e.~position 1) from this
analysis, as they were not part of the temporal order memory decision
(within-event temporal order judgment was on trials 2 and 6). In all
three bands, we observed a significant linear trend for events where
order was later correct {[}delta: F(1,17)=19.88, p\textless{}.001;
theta: F(1,17)=18.82, p\textless{}.001; beta: F(1,17)=29.83,
p\textless{}.001{]}, but not for incorrect events {[}delta:
F(1,17)=.294, p\textgreater{}.1; theta: F(1,17)=1.53, p\textgreater{}.1;
beta: F(1,17)=1.95, p\textgreater{}.1{]}. For beta (but not delta or
theta; p's\textgreater{}.2), there was also a significant position by
memory interaction {[}F(4,68)=3.11, p=.02{]} and a trending interaction
of the linear contrast {[}F(4,68)=3.76, p=.069{]}. Since differences in
trial numbers between conditions can influence estimates of oscillatory
power, we performed an additional sub-sampling analysis where for each
iteration (n=1000) we randomly selected a subsample of trials to match
the number of trials between conditions. Since every subject had more
trials for the remember condition, this analysis essentially matched the
number of trials in the `hits' bin to the number of trials in the
`misses' bin. We performed this analysis for all subjects and then
averaged across the iterations and recomputed our statistics. This
analysis suggested that differences in trial numbers going into the
power estimates did not account for differences in the linear trend as a
function of conditions {[}linear trend for hits: delta: F(1,17)=20.74,
p\textless{}.001; theta: F(1,17)=18.24, p\textless{}.001; beta:
F(1,17)=30.19, p\textless{}.001; linear trend for misses same as
reported above{]}. The significant memory by position interaction
{[}F(1,17)=3.05, p=.02{]} and trending interaction in beta power of the
linear contrast (hits vs.~misses) was also replicated {[}F(1,17)=3.78,
p=.07{]}. In summary, oscillatory power (in delta and theta, but most
robustly in beta) during successful order encoding linearly increased
over trials within an event and this linear pattern was not present when
temporal order later incorrect.

Notably, this linear accumulation effect was not driven by overall
greater power for temporal order correct events. In fact, the opposite
pattern was true: trials for which temporal order was later tested
actually showed greater oscillatory power for misses as compared to hits
{[}delta, 2nd position: t(17)=-2.12, p\textless{}.05; delta, 6 th
position: t(17)=-2.09, p=.051; theta, 2nd position: t(17)=-2.03, p=.057;
theta, 6th position: t(17)=-2.28, p\textless{}.05; beta, 2nd position:
t(17)=-2.26, p\textless{}.05; no other significant differences{]}. Thus,
the pattern of power accumulation across an event was predictive later
order memory, but the absolute magnitude of power (in all three bands)
was predictive of later mnemonic errors.

\subsubsection{Leveraging across-subject variability in power
accumulation model
fits}\label{leveraging-across-subject-variability-in-power-accumulation-model-fits}

Next, we tested the prediction that variability in the extent to which
individual participants fit the accumulation model would be related to
within-event temporal order memory. Specifically, we predicted that
subjects whose oscillatory power dynamics better fit the accumulation
model would show better within-event temporal order memory as compared
to subjects that poorly fit the model. We limited our analyses to the
delta, theta and beta bands averaging across the sensors that
significantly fit the accumulation model (at the group-level).
Consistent with our prediction, beta accumulation model fits were
positively correlated with within-event order memory {[}Figure 4; robust
regression: t(16)=1.73, p\textless{}.05, one-tailed{]}. There was no
such relationship in delta or theta {[}p's\textgreater{}.1{]}. It's
worthwhile to note that if when we analyzed the data with a stricter
statistical threshold (that is, averaging across sensors that exceeded a
p\textless{}.0001 threshold), the correlation became stronger {[}robust
regression: {[}t(16)=2.16, p\textless{}.05{]}. In summary, participants
whose beta power dynamics better fit the event accumulation model showed
relatively better temporal order memory.

\subsubsection{Boundary-related decreases in power predict
across-boundary order memory
cost}\label{boundary-related-decreases-in-power-predict-across-boundary-order-memory-cost}

The previous analyses were aimed to identify oscillatory activity that
followed a linearly increasing pattern within an event. In this next set
of analyses, we test the hypothesis that event boundaries result in a
significant drop in power. Although we've already provided evidence that
oscillatory activity significantly fits our within-event accumulation
model, it remains possible that this is pattern is not driven by a
significant drop in power at boundaries. For instance, a linear trend
(within an event) could exist in the data without any significant
differences between positions. To test this, we extracted delta, theta
and beta power data from sensors showing significant accumulation model
fits (averaging across sensor) for trials in position 6 (i.e.~the last
trial before a boundary) and trials in position 1 (the first trial after
a boundary). We hypothesized that power would be significantly higher
for trials in the last position of an event compared to trials in the
first position of an event. We found this to be true for all three bands
{[}delta: t(17)=4.85, p\textless{}.001; theta: t(17)=5.05,
p\textless{}.001; beta: t(17)=7.52, p\textless{}.001{]}. While this
analysis is not entirely independent from within-event power
accumulation analyses described earlier, it provides corroborating
evidence that event boundaries result in a drop in a significant
oscillatory power.

Next, we tested the hypothesis that boundary-related decreases in power
would influence temporal order memory organization. Our logic was the
following: if power decreases at event boundaries represent flushing of
event content from working memory, then subjects who displayed larger
boundary-related decrease in oscillatory power might also show larger
boundary-related cost in temporal order memory. For our three bands of
interest, we extracted power from sensors showing significant
accumulation model fit from trials in position 6 (i.e.~the last trial in
an event) and position 1 (i.e.~the first trial after a boundary) and
computed a difference score (position 6 minus position 1 power). We
hypothesized that this difference score might represent an individual's
neural sensitivity to event boundaries. Then, for each subject, we
computed a boundary-related order memory cost index by subtracting
overall across-boundary order memory from within-event order memory.
This difference score served as a proxy for how much the boundaries
reduced across-event temporal order memory (relative to within-event
order memory). We removed one outlier whose behavioral difference score
was 2.37 standard deviations from the mean. We found that the
boundary-related decrease in beta power (but not theta or delta,) was
correlated with the magnitude of the across-boundary order memory cost
{[}Figure 5; robust regression: t(16)=2.99, p\textless{}.01{]}. This
correlation remained significant when including the outlier subject in
the analysis {[}robust regression: t(17)=2.58, p\textless{}.05{]}. We
did not observe this pattern in delta or theta (both
p's\textgreater{}.4). When considering the within-event and across-event
trials alone, there was trending effect for within-event {[}robust
regression: t(17)=1.66, p=.06{]} and no effect for across-event
{[}robust regression: t(17)=.83, p=.11{]}. Thus, the subtraction likely
acted as a baseline to remove gross across-subject memory differences
that might have dampened this relationship. This finding suggests that
the magnitude of drop in beta power at event boundaries is related to
the impact of event boundaries on temporal order memory.

\section{Discussion}\label{discussion}

In this study, we recorded oscillatory brain activity during the
encoding of temporally extended experiences and found that temporal
dynamics in delta, theta and beta power tracked the event structure in
our experiment. Specifically, oscillatory power accumulated within an
event, or a sequence of 6 adjacent trials with a shared background
color, and then dropped sharply at boundaries between events. Moreover,
linear increases in power within each event were related to temporal
order memory success, suggesting that these frequency bands (delta,
theta and beta) may play an important role in the organization of
temporally extended memories. These findings contribute to a growing
body of evidence suggesting that neural oscillations play a prominent
role in the formation of episodic memories. Furthermore, our unique
across-trial regression approach to the MEG data highlights the
importance of considering oscillatory dynamics across longer time scales
than what has been typical in the literature.

\subsubsection{Delta/Theta oscillations and temporal memory
formation}\label{deltatheta-oscillations-and-temporal-memory-formation}

The pattern of accumulating theta power within an event is consistent
with prior research demonstrating that increases in theta power are
correlated with working memory load (Jensen and Tesche). Furthermore, a
recent EEG study found that (frontal) theta power was greater during the
maintenance of temporal order information as compared to item
maintenance alone (Hsieh and Ranganath). In the current study, subjects
may be accumulating/maintaining within-event item representations as an
event unfolds. If that is the case, then one might expect theta power to
increase with each additional maintained item (as working memory load
increases). This is exactly the pattern of results that we observed.
Moreover, theta power accumulated specifically when the order of the
items within the event was successfully encoded (and not when temporal
order was incorrect), suggesting that dynamic increases in theta power
across an event may support temporal order encoding. One interpretation
of this pattern is that accumulating theta power within an event is a
by-product of neural processes supporting the maintenance of multiple
items in working memory. Neural representations of items that are
co-active in working memory might become associatively bound to one
another, which could support the recovery of temporal information on a
later test, and explain why accumulating theta power is related to
temporal order memory success. This interaction between a working memory
`buffer' and long-term associative memory is predicted by memory models
that posit that items may become associatively bound in long-term memory
via a `buffer'-like process active during encoding (Lehman and
Malmberg).

There is a good deal of evidence now that theta oscillations also play a
role in long-term memory encoding (REFS), and some recent work has
argued that theta oscillations may play a more specific role in
associative binding (REFS). Summerfield and Mangels (2005) found that
successful word-color associative memory was accompanied by an increase
in frontal theta power, as well as an increase in frontal-posterior
theta coherence. In another relevant study, Staudigl and Hanslmayr
(2013) recorded oscillatory activity while subjects encoded words that
were superimposed on a background movie. Then, they test memory for the
words, where some trials were presented intact (same word, same movie)
while other trials were rearranged (same word, different movie).
Interestingly, for intact test trials they found that greater theta
power during encoding was predictive of item memory. However, for
rearranged test trials, the reverse pattern (i.e.~relatively less theta
power during encoding) was predictive item memory. This suggests that
theta power during encoding indexes item-context binding, since
presumably greater item-context binding would facilitate retrieval in
the intact condition, but not in the rearranged condition. Taken
together, these results are consistent with the idea that theta
oscillations may be an index of neural processes generally supporting
associative binding, whether its within-trial (i.e.~item-context) or
across multiple trials (i.e.~item-item).

Our source analyses suggest that delta and theta power accumulation
appear to be localized to the left hippocampus and right medial temporal
lobes (MTL), respectively. This is consistent with previous research
that has shown that theta oscillations in the MTL appear to play a
critical role in successful memory encoding (CITATIONS). While some
caution is warranted due to the difficult problem of deep source
localization (REFS), given the robust literature on MTL contributions to
memory encoding paired with the localization results here, we
hypothesize that the within-event delta/theta power accumulation effects
we observed in this study are indeed generated by MTL structures. Future
studies could utilize convergent techniques with higher spatial
resolution (i.e.~intracranial EEG) to validate this localization.

\subsubsection{The Role of Beta Oscillations in the Maintenance of the
Status
Quo}\label{the-role-of-beta-oscillations-in-the-maintenance-of-the-status-quo}

In addition to accumulating theta power, we also found a robust
relationship between accumulating beta power and temporal order memory
within events (Figure XX). Like theta power, accumulating beta power
within an event was predictive of temporal order memory success.
Furthermore, subjects who showed larger boundary-related drops in beta
power also had better within- vs.~across-event temporal order memory,
suggesting that beta power dynamics have bearing on memory organization.

In a review of the literature, Engel and Fries (2010) propose that beta
oscillations might index the `maintenance of the status quo'. This
theory was primarily inspired by research in the motor domain which has
shown that beta power is elevated during sustained motor operations
(i.e.~holding a posture) and decreases during voluntary movements
(REFS). The authors suggest that if beta oscillations index the
maintenance of a particular motor state, then they may also be elevated
during the maintenance of a cognitive set (and in other situations of
high levels of endogenous attention) and decrease in the face of novel
or unexpected events. Our results fit well within this framework: we
found that beta oscillations increased monotonically as the event
unfolded, and rapidly dropped at event boundaries (i.e.~when the
background changed to a new novel color). Bolstered by this theoretical
work, we speculate that the beta oscillations we observed here might
index the maintenance of a `cognitive set', or perhaps an `event model',
a construct that is central to event segmentation theory (REFS).

Beta oscillations have also been directly linked with successful memory
encoding. In a series of intriguing studies, Hanslmayr and colleagues
found that beta desynchronization in the frontal cortex is predictive of
subsequent memory and is also positively correlated with BOLD activity.
Furthermore, the authors found that entrainment of the frontal cortex to
the beta frequency (but not theta or alpha entrainment) during encoding
selectively impaired later memory (Hanslmayr). A possible explanation
for this finding is that beta desynchronization may actually lead to a
richer, more unique neural representation (REFS). This prediction stems
from the idea that theoretically, desynchronized brain activity has a
higher information capacity than synchronized activity (REFs).
Connecting this idea to the current work, we speculate that accumulating
beta power (i.e.~increases in synchronized activity) may represent the
maintenance of a prior cognitive set (or event model). At event
boundaries, when the background is updated to a novel color, the sharp
drop in beta power may reflect bottom-up processing of novel perceptual
information, which may require relatively greater processing bandwidth.
Thus, while beta power accumulation may be beneficial to forming
associative relationships among a group of items, at the same time
greater beta power may be detrimental to encoding in scenarios where
item distinctiveness is important for memory. Future work should
directly test this tentative hypothesis.

While low frequency power accumulation was localized to MTL structures,
beta power accumulation was localized to regions including the frontal
and parietal cortex (as well as the MTL; see Supplemental Table XX full
list). Prior studies have found that synchronized beta activity between
frontal and parietal regions increases under conditions of high top-down
endogenous attention (Buchsman and Miller; Buchsman and Miller). Thus,
the observed increase in beta power within an event could reflect an
increase in exogenous attention necessary to maintain accumulating items
in working memory. Thus, at boundaries, the drop in beta power could
reflect a shift from endogenous to exogenous attention.

Finally, we observed overlap in the source localization of theta and
beta in the medial temporal lobe (as well as the insula). Both theta
(CITATIONS) and beta (Kopell, Traub) have been observed in medial
temporal lobe structures. Empirical data and computational models
support the idea that while higher frequency bands (such as gamma) can
be used for relatively local neural computations, it's possible for
distal brain regions to synchronize at lower frequencies
(\textless{}gamma). Thus, both of beta and theta in the hippocampus
could be involved in the integration of representations that are
distributed across different brain regions in the service of episodic
memory encoding. Future work should address if and how these frequencies
might interact.

\subsubsection{Summary}\label{summary}

In summary, these findings suggest that accumulating oscillatory power
across multiple trials in the delta, theta and beta band play a role in
memory encoding and organization.

\begin{figure}[htbp]
\centering
\includegraphics{figures/figure1}
\caption{This is the caption}
\end{figure}
